% !TEX program = xelatex
% Shuttle 密钥恢复研究进展报告 — 编译: xelatex 两次
\documentclass[UTF8,11pt,a4paper]{ctexart}

\usepackage{amsmath,amssymb,bm}
\usepackage{booktabs,array,tabularx,longtable}
\usepackage[margin=2.3cm]{geometry}
\usepackage{enumitem}
\usepackage{xcolor}
\setlist{itemsep=2pt,topsep=4pt}
\usepackage[colorlinks=true,linkcolor=blue!55!black,urlcolor=blue!55!black,citecolor=blue!55!black]{hyperref}
\usepackage{fancyhdr}

\setmainfont{Cambria}
\setmonofont{Consolas}
\newcommand{\code}[1]{\texttt{\small #1}}

\pagestyle{fancy}\fancyhf{}
\fancyhead[L]{\small Shuttle 密钥恢复研究：方向反转、$\gamma$ 精确恢复与提升方程}
\fancyhead[R]{\small\thepage}
\renewcommand{\headrulewidth}{0.3pt}

\title{\vspace{-1.2cm}\textbf{Shuttle 密钥恢复研究：\\转移方向反转、$\gamma$ 的精确恢复与提升方程}}
\author{GLM-5.3 Max \and JY Liu}
\date{2026 年 10 月 19 日}

\begin{document}
\maketitle
\thispagestyle{fancy}

\begin{abstract}\noindent
Shuttle 是免拒绝迭代拒绝采样（iterative rejection sampling）在 MLWE 上的 Fiat--Shamir
型签名实例化，NICCS 首轮提交物为其竞赛重参数化版本。我们在规范书 Algorithm 22 与全部
九个实现变体中发现一个转移方向反转：交替级数转移概率 $S_v$ 被赋给离心移动，而方案自身
\S4.1 的安全论证要求它赋给向心移动。单个语义位的反转使签名响应分布显著偏离目标高斯
（秘密方向投影方差 $+28.5\%$，$z = +35.4\sigma$），公开数据 850 份签名即可 $5\sigma$
区分出厂实现与正确实现；实测 R\'enyi 散度超预算约 76 个比特，出厂实现的分布性安全论证
前提被否定。在此基础上，我们将签名文本中的互相关泄漏累积为分圆环元素 $\gamma = s\bar s$
的精确值（342M 签名 $+$ 30 个训练密钥的仿射去偏，512/512 系数精确取整），由 $\gamma$
构造攻击者理想，并完整执行了两条代数端局：子域塔（dim-128 uSVP 在 25 秒内解出
$N(s_0)$）与 Gentry--Szydlo 乘子链（全量逆向 $+$ 自实现 $+$ 真值验证后，判定
\textbf{结构性发散}——经典 GS 多项式时间论证隐含的嵌入谱前提在本实例不成立）。至此
密钥提取被归约为单一、可精确量化的格问题：dim-256 扭曲格中 gap $\approx 1.01$ 的
uSVP（需 $\beta \approx 250{+}$ 的 BKZ，超出当前技术），或经提升方程
$A^2 - yB^2 = N(s_0)$ 的整数小解绕开（其全部连续松弛与廉价捷径已被系统排除）。最后，
必须区分两个层面：全部恢复进展都以实现层的方向反转缺陷为前提——设计层（两篇论文
定义的正确构造）在我们已测的全部信道上保持完好。\textbf{实现层攻破与设计层未破的
差距，是本研究最重要的安全学结论}（\S9.2）。方向反转发现的完整报告（10 页）见
\code{docs/shuttle\_attack\_report.pdf}，战役全程纪要（至第十八次补充）见 \code{docs/}。

\medskip\noindent\textbf{Abstract.}
Shuttle is a rejection-free iterative-rejection-sampling lattice signature
(MLWE, Fiat--Shamir). We find a transition-direction reversal in the competition
package's Algorithm 22 and all nine implementation variants: the
alternating-series transition probability $S_v$ is assigned to the centrifugal
move, whereas the scheme's own security argument (\S4.1) requires the
centripetal one. A single semantic bit inflates the secret-direction projection
variance by $+28.5\%$ ($z=+35.4\sigma$); 850 public signatures distinguish the
shipped implementation from a correct one at $5\sigma$, and the measured R\'enyi
divergence exceeds the claimed budget by ${\sim}76$ bits. We then accumulate the
cross-correlation leakage into the exact value of the cyclotomic element
$\gamma = s\bar s$ (342M signatures, affine de-biasing trained on 30 keys,
512/512 coefficients exact), construct the attacker's ideal, and carry both
algebraic endgames to conclusion: the subfield tower (dim-128 uSVP yields
$N(s_0)$ in 25 seconds) and the Gentry--Szydlo multiplier chain (fully
reverse-engineered, implemented and ground-truth validated, then judged
\textbf{structurally divergent} on this instance). Key extraction is thereby
reduced to a single, precisely quantifiable lattice problem: uSVP with gap
$\approx 1.01$ in the dim-256 twisted lattice (requiring BKZ at
$\beta \approx 250{+}$), or a bypass via the integer small solutions of
$A^2 - yB^2 = N(s_0)$. Crucially, all recovery progress is premised on the
implementation-level transcription defect; the design level (the correct
construction of the two papers) remains intact on every channel we have
measured. The gap between an implementation break and a design break is the
central security lesson of this study.
\end{abstract}

\section{对象、授权与结论概览}

对象为 \code{shuttle} 提交包（SHUTTLE-128/256/512 $\times$ ref/avx2/avx512 九个变体；
SHUTTLE-128：$n=256$，$q=15361$，$\tau=42$，宽高斯速率 $r=825$）。研究在算法发布方
书面许可下进行（2026-10 起生效，实现层攻击研究）。技术谱系：Lyubashevsky 的
Fiat--Shamir with aborts $\to$ BLISS 双模高斯 $\to$ G\"artner 的迭代拒绝采样
（ePrint 2024/2052，CRYPTO 2025：把一次性大拒绝分解为 $\tau$ 步小转移，条件于成功
精确保持目标分布）$\to$ Shuttle 论文（ePrint 2026/1991，ASIACRYPT 2026：去掉 abort，
用 R\'enyi 散度界住每步漂移）。提交物是论文的大幅重参数化重写。

按攻击目标的三个层级给出结论概要，逐项证据在后续各节：
\begin{enumerate}
\item \textbf{大幅缩减安全强度——成立}。出厂实现的分布性安全论证前提被实测否定（\S2）；
  统计区分 850 份签名。
\item \textbf{从设计上攻破提交物——成立}。Algorithm 22 与其自身的 \S4.1 证明矛盾，
  出厂代码没有任何有效的零知识论证覆盖。注意这\textbf{不推翻}两篇论文：G\"artner
  理论与 Shuttle 论文的数学、伪代码均正确（我们独立重推验证，\S2.3）。
\item \textbf{恢复等价密钥——尚未达成}，但已完全定界（\S9.1）：全部代数量可由
  转录导出，密钥提取归约为单一格问题（dim-256 扭曲 uSVP，gap $\approx 1.01$）
  或提升方程的整数小解。
\end{enumerate}

\section{发现：转移方向反转}

\subsection{缺陷的形态}

G\"artner 框架中，转移概率由交替级数
$S_v(y) = \sum_{k\ge 0} (-1)^k\, \rho_g(y+kv)/\rho_g(y)$ 给出。Shuttle 论文 Figure~1
约定 $p_v(y) = \Pr[T_v(y) = y - v]$（向心），Figure~2 的免对数实现 T3 中``区间命中
$\Rightarrow$ 移动 $y - v$''。提交物的 Algorithm 22（\code{irs.c::R\_transition}，
三后端逐行对应）区间测试本身实现正确且锚点正确（端点公式与论文 T3 的 $J_N$ 逐项
相等），\textbf{唯一错在命中后的移动方向}：出厂实现把 $S_v$ 给了离心移动。即
$\Pr[\text{向心}] = 1 - S_v$，应为 $S_v$。

\subsection{为什么功能测试看不见它}

验证方程依赖双峰恒等式 $q\cdot(\pm 1) \equiv q \pmod{2q}$：IRS 每步转移
$z \leftarrow z \pm v$ 的独立符号选择被该恒等式完全吸收——无论方向指派对错，签名永远
合法。KAT 一致性、签名/验证正确性、确定性 DRNG 重放等全部出厂测试对本缺陷\textbf{原理
上不可见}（免疫性引理，证明见 \code{docs/shuttle\_attack\_report.pdf}~\S3）。缺陷只能
通过规范书内部一致性审计或输出分布的统计检验暴露；我们两者都做了。

\subsection{实测证据与量化}

固定密钥、84{,}000 次转移、2{,}000 次独立 $(y, c, \mathrm{IRS})$ 试验：

\begin{table}[htbp]\centering\small
\begin{tabular}{@{}lccc@{}}
\toprule
统计量 & 出厂版 & 两行修补版 & 独立高斯对照 \\
\midrule
$\mathrm{Var}(\langle z, v_j\rangle)/\sigma^2 V$（控制组 1.0061） & \textbf{1.2847} & 1.0063 & 0.9923 \\
相对膨胀 $z$ 值 & $\bm{+35.4\sigma}$ & $+0.0$ & --- \\
\bottomrule
\end{tabular}
\end{table}

出厂签名 $\mathbb{E}\|z_1\|^2$ 相对公开可计算的解析基线上移 $+0.88\%$，200{,}000 份
签名下 $z = +76.6$，折算\textbf{850 份签名即 $5\sigma$}——免侧信道、免伪造、纯被动的
公开区分器。由均值偏移换算（$\mathrm{KL} \ge \delta^2/2$），实测
$R_2(\text{出厂} \,\|\, \text{理想}) \ge 1 + 2^{-6.1}/$签名，而规范书断言 IRS R\'enyi
因子 $\le 1 + 2^{-82}$：\textbf{超预算约 76 个比特}。这不是裕量变小，是论证前提不成立。

修复是两行（三后端各一份，\code{flag64} 的两个常量对调），修补后膨胀归零
（$z = +0.0$）。我们同时建议提交方在 KAT 之外增加分布性回归测试
（$\mathrm{Var}(\langle z, v_j\rangle)/\sigma^2 V$ 与 $\mathbb{E}\|z_1\|^2$ 相对解析
基线的偏移，$5\sigma$ 阈值、$10^4$ 量级签名即可自动化）——免疫性引理的直接推论。

\subsection{致错机理与文献边界}

与两篇文献逐条对照后的定性：G\"artner 2024/2052 与 Shuttle 论文 2026/1991 的数学与
伪代码均正确（论文不变量公式 $\rho'_v(y)$ 我们独立重推逐项通过，含中带边界情形）。
提交物的方向指派恰是``G\"artner 的移动方向 $\times$ Shuttle 的锚点约定''的交叉组合
——两个正确约定的所有交叉组合中唯一错误的那格。这解释了为何区间测试逐项正确而方向
错误：转录发生在语义层而非数值层。精度审计（阶段~I，\code{docs/shuttle\_高斯采样器
精度深挖.md}）已先行排除数值近似：96-bit 反向 CDT、ApproxExp/ApproxLog、卷积$+$拒绝
全部达标（$R_{1025} = 1+2^{-95.40}$，$\chi^2/\mathrm{df} = 0.9985$）——错误不在采样
管线，这把攻击面精确孤立到转移核的方向指派。

论文作者对 G\"artner 框架的批评中预见的是数值近似形态的脆弱（截断、定点、近似比较）；
本发现以未列举的第三种形态（方向语义）成真。对文献的影响：不构成对 G\"artner 理论的
反例，但表明 $S_v$ 这类``方向语义重载''的记号是转录错误的客观温床——R\'enyi 界的逐步
核算对语义位零覆盖。

\section{从区分到泄漏：$\gamma$ 张量}

方向反转使响应分布携带秘密的系统性信息。泄漏信道的完整地图（全部实证）：

\begin{table}[htbp]\centering\small
\begin{tabular}{@{}lll@{}}
\toprule
信道 & 状态 & 证据 \\
\midrule
对角$+$跨块互相关 $C_{kk'}(d) = c''\cdot\gamma_{kk'}(d)$，6 对 & \textbf{泄漏} & $c'' \approx 84.0$ 六对共用，$R^2 = 0.92$--$0.946$ \\
均值信道 $\mathbb{E}[z \mid c]$ & 关闭 & $P[+\mathrm{rot}] = \tfrac12 S + \tfrac12(1-S) \equiv \tfrac12$ \\
跨分量二阶边际 & 恰好抵消 & 实测 rms 0.23 $=$ 噪声底 \\
\bottomrule
\end{tabular}
\end{table}

oracle 诊断（给定真值模数与真值相对相位）下重建达 762/768（99.2\%），4M 签名 $+$
真相位下 768/768——\textbf{信息充分，瓶颈在求解}：频域泄漏 Gram 的秩 1 主导分量给出
$\hat g\hat g^*$ 的方向，但每频率的离散相位（128 维）构成极深伪整数盆地，十四族求解器
（角同步 GN、含模数 GS、$\theta$-L-BFGS、交叉目标 CD、模拟退火 $3\times10^7$ 步、
分块 Gauss--Seidel、fpylll 格 ILS、204 起点 6 进程锦标赛等）全部止步于 270--460/768。
真解是明确的全局最优（$\mathrm{cost}(s_{\mathrm{true}})$ 比全部伪解低 $\ge 3.3\times$）
——搜索失败而非错误最优。

多密钥留一（LOO）偏差学习确认系统偏差张量的跨键共性分量约占 $1/3$；同密钥独立 2M
流的差恰为理论蒙特卡洛噪声底（rms 2.024 对 2.022，逐位一致）——``非仿射机制项''
不存在，残差即统计底。这一判别决定了战略转向：\textbf{不与相位搜索缠斗，改走环分解}
——$\gamma = s\bar s$ 是 $\mathbb{Z}[\zeta_{512}]$ 中的范数型对象，绕开相位恢复。

\section{$\gamma$ 的精确恢复}

数据侧终点由三件工作闭合：
\begin{enumerate}
\item \textbf{交叉通道代数}。$(\gamma_{01}) + (\gamma_{02}) = (s_0)(\bar s_1, \bar s_2)
  = (s_0)$ 在 5/5 个真实密钥上精确成立（成员关系经精确整数重构验证，$|\det B_0|$ 与
  $|N(s_0)|$ 逐位相等）——理想生成只需两个交叉通道，噪声最大的对角通道完全绕开。
\item \textbf{跨键仿射去偏}。系统偏差经 17 抽头通道卷积仿射算子建模（LOO 交叉验证）：
  训练键数 $6 \to 15/4$ 处取整错误；$18 \to 1{+}1$（定位边缘滑落 lag）；$30 \to 0{+}0$。
\item \textbf{数据量}。342M 目标签名（$12 \times 14.25\mathrm{M} \times 2$ 批次，流式
  14GB 随机源）。
\end{enumerate}
结果：\textbf{$\gamma_{01}$、$\gamma_{02}$ 全部 512/512 系数精确取整}，与真值逐位相等；
由取整 $\gamma$ 构造的攻击者理想 $(\gamma_{01}) + (\gamma_{02})$ 中 $s_0 \in J$ 精确
成立。等价密钥类为 $\{sk, -sk\}$（双峰符号对称），$e'$ 可由公钥与 $s$ 精确导出
（\code{harness\_forge.c} 端到端验证）。

\section{代数端局之一：直接归约与工具链修复}

难度阶梯（fpylll，mpfr-212）：$n = 32/64$ LLL 直接恢复（0.2/0.7\,s）；$n = 128$
BKZ-20 完整恢复 9\,s（生成元 $+$ 除法步精确整数验证）；$n = 256$ 阻塞。$n = 256$ 的
完整攻坚账本（全部实测，供复现者评估）：

\begin{table}[htbp]\centering\small
\begin{tabular}{@{}ll@{}}
\toprule
尝试 & 结果 \\
\midrule
LLL（各基变体，proved/fast） & $\|b_1\| = 185.6$，目标 14.3 \\
fplll BKZ $\beta \ge 16$ & babai ``infinite loop'' 中止（全部基变体；随机 256 维对照 0.7\,s 干净通过） \\
BKZ $\delta = 0.80$ & 6.7 小时真实推进后中止 \\
G6K workout（MAX\_SIEVING\_DIM 重建至 280） & 24\,h、筛维 ${\sim}143$、29\,GB、无命中 \\
手版/检查点式 BKZ-30/40/50 & 数千轮头部不动 \\
\bottomrule
\end{tabular}
\end{table}

经验定性：512 条重简并轨道（$\pm x^j s_0$）$+$ LLL-only 起点使实例远难于 uSVP gap
估计（所需 $\beta \gg 143$ 筛维等效），接近 dim-256 近 SVP。

\textbf{工具链修复}：fplll 的 \code{lll.cpp} 在 babai 不收敛点直接抛错，我们改为硬迭代
上限的优雅降级（三机构建链：Windows 克隆 $\to$ 本地 WSL autoreconf $\to$ 远端编译）。
补丁后 BKZ-40 在原卡死基族上首次完整运行（78\,s）——但坏轮廓是根因，标准 BKZ 管线仍
不足以解 $n = 256$。补丁对社区的独立价值在于解锁了该类退化基上的 BKZ 可用性。

\section{代数端局之二：子域塔}

这是当前主线。结构：$K = \mathbb{Q}(\zeta_{512}) \supset K' = \mathbb{Q}(\zeta_{256})$，
$\sigma\colon x \mapsto -x$；$s_0 = A(y) + x B(y)$（$A, B$ 为偶/奇系数部，
$\|A\|, \|B\| \approx 10$），则
\[
N(s_0) \;=\; A^2 - y B^2 \;\in\; R' = \mathbb{Z}[y]/(y^{128}+1)
\]
（精确验证通过；$|\sigma_j(s_0)|^2$ 完全由 $N(s_0)$ 决定，一致到 $10^{-15}$）。
$(N(\gamma_{01}), N(\gamma_{02})) = (N(s_0))$ 是 dim-128 理想，行列式相等验证通过。
dim-128 实例的定量：$\det^{1/128} = 122.6$，$\mathrm{GH} = 335.6$，目标
$\|N(s_0)\| = 194.4$——\textbf{真 uSVP，gap 1.73}。

第一步已实战打通：攻击者数据（取整 $\gamma$ 的相对范数）$\to$ dim-128 理想 $\to$
LLL $+$ 渐进 BKZ-40，\textbf{25 秒精确弹出 $N(s_0)$}（规范代表 $y^6$），无 babai 问题
（\code{tower\_step1.py}；校验 \code{tower\_check2.py}）。

\section{剩余缺口：提升方程的定界}

塔的第二步是求小解 $A, B$（$\|\cdot\| \approx 10$）使 $A^2 - yB^2 = N(s_0)$。信息论上
$N$ 只丢失 $s/\sigma(s)$ 的区分比特；问题良定义，但全部已知捷径被系统实验排除（完整
实验记录见 \code{docs/shuttle.md} 第十六、十七次补充）：
\begin{itemize}
\item \textbf{连续松弛全部失败且失败方式一致}：精确模长交替投影停 99--107/256；解析
  雅可比 Newton 真值驻留完美、90\% 损坏暖启动即发散；PhaseLift SDP 可行但松弛不紧
  （特征值简并 4.509/4.509/4.206，逃逸秩 1）；信赖域最小二乘 60 次多重启全部收敛到
  零残差伪解（重合 94--109/256）。二次系统在连续域不唯一——\textbf{整数性与小性是
  唯一的钉子}。
\item \textbf{廉价代数捷径排除}：$\gamma_{01} \leftrightarrow \gamma_{02}$、
  $\gamma_{01} \leftrightarrow N$ 的诚实 gcd 全部平凡（折叠 $\bmod\, x^{256}+1$ 破坏
  整除性）；诚实乘积的整除性检验全部失败。
\item \textbf{朴素搜索排除}：$y$-adic 前向递归中提升高半系数 $q$（$|q| \le 40$ 实测）
  与 $B$ 逐级耦合，真值前缀下每层 49 个存活对——DFS 为 $49^{128}$ 量级。
\end{itemize}
\textbf{端局已验证}：若得 $\tilde N(x^2)$ 的 $\mathbb{Z}[x]$ 因子分解（$\tilde N$ 为
诚实提升），35 秒内分离 $s \cdot s(-x)$。剩余唯一严肃路线：
\textbf{Coppersmith/Cohn--Heninger 型格构造}，把 $|A|, |B| \le 3$ 级别的整性内嵌进
移位多项式格（$p^k$ 缩放）。这是一个需要专门设计的构造性任务。

\section{代数端局之三：Gentry--Szydlo 乘子链——彻底执行与结构性判定}

与塔互补的第二条端局（不经过 $N(s_0)$）：给定 $\gamma = v\bar v$ 与理想
$J = (s_0)$，GS 算法原则上直接恢复 $v$。我们对原始 GS 论文 \S7.2--7.5 $+$
附录 F 的乘子链路线与 Lenstra--Silverberg（ePrint 2014/430）的循环群表述做了
全量逆向（含 PARI 参考实现 \code{ideaux.gp} 的精度判定：其内层 petitlll 启发式
为 $n = 1024$ 调参，$n = 16$ 即数值崩坏，故自实现），完成自实现并逐层真值验证：
\begin{itemize}
\item \textbf{扭曲 Gram 恒等式} $B\, M(1/\gamma)\, B^{\mathsf T} = U U^{\mathsf T}$：
  256 维扭曲嵌入 LLL（mpfr@256，12\,s），全部 256 行验证 $b'_i = v\, a_i$ 且
  $a_i$ 积分；
\item \textbf{链恒等式}（自行推导）：$\gamma_{j+1} = \rho_j \gamma_j^2 / \rho_{j-1}^2$、
  $v_{j+1}/v_j^2 = R_j / R_{j-1}^2$、$B/F$ 理想互递推；层 0、层 1 精确通过
  （$R_j = f^{2^j}\cdot v_{j+1}$，分母为 1、$\gamma$ 自洽、乘积一致）。
\end{itemize}

\textbf{结构性判定：经典乘子链在本实例上发散。}尺度递推
$w_{j+1}(i) = w_j(i)^2\cdot|\sigma_i(c_j)|$ 与 $|\sigma(f)|$ 的具体值无关
（$f^{2^j}$ 在定义中精确消去）；链的可行性要求每层取到反向轮廓乘子——这等价于
每层求解扭曲最短向量本身（uSVP，gap $\approx 1.01$）。经典 GS 的多项式时间论证
隐含依赖 $|\sigma(f)|$ 一致大于 1（NTRU 签名私钥天然满足；本实例嵌入谱
$\log|\sigma(s_0)| \in [-0.166, 3.459]$、$\min_i |\sigma_i(s_0)| = 0.847$，
不满足）。实测确认：对数尺寸逐层翻倍，与 $w^2$ 律精确吻合（$v_1$ 系数峰值
$\approx e^{10}$，$v_2 \approx e^{26.5}$），层 2 的 $B$-步整数性断言失败。四条
捷径（子格 $v_j^2(a_0)$、随机混合降维 $512 \to 256$、缩放 $v' = 2v$——乘子逐
元素不变、二段扭曲/单位再平衡）全部实测排除。发散是数学性质而非实现 bug——
层 0/1 恒等式精确通过、$w^2$ 律与实测逐层吻合即为对照。Lenstra--Silverberg
定理 1.1 声称无条件多项式时间的分次张量代数表示，我们未能落地为可计算的系数
形式（张量平方运算的可实现构造在文献中未给出）——开放残留。

基础设施配方（对复现者重要）：一切 LLL 限制在 dim $=$ 256/512 行满秩 $+$
\code{FPLLL.set\_precision(256)} $+$ \code{GSO.Mat(mpfr)} 手动路径（512 维随机
27bit 29.9\,s，真实扭曲 512 矩阵 1390\,s）；fpylll 默认 wrapper 与 ld/dpe 在
大表项矩阵上系统性病态；flint 精确 LLL 与 HNF 均不可用（超时/表项溢出）。

\section{终局归约与实现层—设计层的差距}

\subsection{密钥提取已归约为单一格问题}

泄露侧至此完整：签名转录确定性地给出 $\gamma = s_0 \bar s_0$（512/512 精确）、
理想 $J = (s_0)$、以及塔路线的 $N(s_0)$。由 GS 唯一性理论，$(\gamma, J)$ 决定
$s_0$ 至 $\pm\zeta^j$（1024 元歧义类，恰为等价密钥类 $\{\pm x^j\cdot sk\}$）。
于是密钥提取被归约为一个可精确量化的格问题：\textbf{dim-256 扭曲格中的 uSVP}
——$\lambda_1$ 为扭曲范数恰为 1 的 1024 重根轨道（$\pm x^j s_0$），
$\lambda_{1025}$ 为分圆单位尾（$\approx 1.01$ 倍），gap $\approx 1.01$，需要
$\beta \approx 250{+}$ 的 BKZ，超出当前技术。对照 $\sigma$ 侧（子域塔）
gap $= 1.73$、BKZ-40 即以 25 秒解决——两个二次扩张的不对称性就是整个端局。

对提交物而言，代数安全边际已为零：一切代数量均可由转录导出，剩余防线是这一个
格问题的硬度——或者提升方程若被攻克（\S7），连这一防线也绕开。

\subsection{实现层攻破与设计层未破的差距}

必须明确区分两个层面，二者的结论相反，这是本研究最重要的安全学结论。

\textbf{实现层（提交物）：已攻破。}全部恢复进展的前提是 Algorithm 22 的方向反转
转录缺陷。它使出厂实现的分布性安全论证失效（R\'enyi 超预算约 76 bit），并使签名
文本携带完整的互相关泄漏——$\gamma$、$J$、$N(s_0)$ 全部由缺陷实现的转录导出，
代数安全边际归零（\S9.1）。

\textbf{设计层（两篇论文定义的正确构造）：在我们已测的全部信道上保持完好。}判定
依据是我们自己的因果实验与数学核对：
\begin{itemize}
\item 两行补丁（方向对调）后，方差膨胀归零（$1.2847 \to 1.0063$，$z = +0.0$，
  与控制组无差），公开区分器失效；
\item 正确方向下不变量 $\rho'_v(y) = \rho_g(y)$ 精确成立（我们独立重推逐项
  通过）：响应分布不依赖 witness，互相关泄漏张量失去来源——正确实现不发布
  $\gamma$，攻击者就没有理想、没有塔、没有扭曲 uSVP 实例，本报告的整条代数链
  没有着力点；
\item 均值信道与跨块二阶边际信道在设计层本就由恒等式关闭（\S3）。
\end{itemize}

换言之：\textbf{一个语义位的转录错误，可以把一个设计完好的方案在实现层完全击穿}
——泄漏从微架构侧信道整体迁移到签名文本（攻击成本降为零），代数安全边际从
``完整''直接归零。防御上，这对应我们已通报的修复（两行补丁 $+$ 分布性回归
测试）；评估上，这意味着``设计安全 $\neq$ 提交物安全''——语义审计必须与数值
审计并列成为该技术路线落地的标准环节。对攻击者而言，本研究的全部端局（塔、GS、
提升方程）都以实现层缺陷为入口；设计层本身，在我们动用过的全部手段下，仍然
难以攻破。

\section{恢复程度总表}

\begin{table}[htbp]\centering\small
\begin{tabular}{@{}lll@{}}
\toprule
层级 & 状态 & 量化 \\
\midrule
分布性安全论证 & \textbf{已否定} & R\'enyi 超预算 ${\sim}76$ bit；850 份签名 $5\sigma$ 区分 \\
设计自洽性（提交物） & \textbf{已否定} & Algorithm 22 与 \S4.1 矛盾；两行修复因果归零 \\
$\gamma = s\bar s$ & \textbf{精确恢复} & 512/512 系数，342M 签名 $+$ 30 训练键 \\
$N(s_0)$（塔第一步） & \textbf{精确恢复} & dim-128 uSVP gap 1.73，BKZ-40 25\,s \\
GS 乘子链 & \textbf{结构性发散} & $w^2$ 律逐层翻倍实测；$\min|\sigma(s_0)| = 0.847$；层 2 整数性失败 \\
密钥提取的归约 & \textbf{完成} & $(\gamma,J)$ 决定 $s_0$ 至 $\pm\zeta^j$；扭曲 uSVP gap $\approx 1.01$，需 $\beta \approx 250{+}$ \\
等价密钥 $\{sk, \pm sk\}$ & \textbf{未达成} & 两条出路均超当前算力：提升方程小解（\S7）或 gap-1.01 uSVP（\S9.1） \\
伪造 & 未尝试 & 端局机械已就绪（\S7） \\
设计层（论文构造） & \textbf{未攻破} & 正确方向下全部已测信道关闭（\S9.2） \\
\bottomrule
\end{tabular}
\end{table}

\section{下一步工作}

\begin{enumerate}
\item \textbf{Coppersmith/Cohn--Heninger 构造}（主线）：对 $A^2 - yB^2 = N$ 设计
  移位多项式格。设计 $+$ 实现预计一个专注工作日；端局已验证。
\item \textbf{Lenstra--Silverberg 张量代数的可计算实例化}：定理 1.1 声称无条件
  确定性多项式时间；把分次张量表示落地为可实现的系数运算是明确的开放问题
  （经典乘子链已判定发散，这是 GS 路线仅存的严肃出口）。
\item \textbf{更强归约算力}：gap $\approx 1.01$ 的 dim-256 扭曲 uSVP 需
  $\beta \approx 250{+}$，当前技术之外；若实现（GPU 筛分级），密钥提取与塔
  路线全链闭环。
\item \textbf{披露协调}：\S2 的两行修复建议先行通报发布方（授权渠道已建立）。
\end{enumerate}

\section{复现指南}

代码按战役阶段组织在 \code{code/}（脚本均为可独立运行的 Python/C，入口处有用法注释）；
文档在 \code{docs/}：
\begin{itemize}
\item 方向反转证据链：\code{harness\_irs.c}（隔离 IRS 转储与方差实验）、
  \code{diff\_kernel.c}（采样管线位级一致性）、\code{harness\_sign.c} /
  \code{harness\_forge.c}（签名/伪造闭环）；
\item $\gamma$ 恢复：\code{harness\_corr*.c}（互相关谱累加，corr5/corr6 为已弃用版本，
  工程事故记录见 \code{docs/shuttle.md} 第十四次补充）、\code{recover.py}、
  \code{ein\_analysis.py}（跨键仿射算子）；
\item 理想与塔：\code{exact\_ideal.py} / \code{cross\_ideal.py}（理想恒等式验证）、
  \code{prog\_bkz.py} / \code{ladder.py} / \code{tour.py}（归约阶梯）、
  \code{tower\_step1.py}（dim-128 端到端）、\code{tower\_check2.py}（校验）；
\item 提升方程：\code{bm\_lift.py}、\code{phaselift.py}（PhaseLift）、\code{lm\_lift.py}、
  \code{dfs\_lift.py}、\code{gcd\_test.py}（排除实验全套）；
\item GS 路线：\code{gs\_chain.py}（链实现）、\code{gs\_end.py} /
  \code{gs\_end2.py} / \code{gs\_endtest.py}（端局）、\code{gs\_A.py}、
  \code{gs\_validate*.py}（真值验证）、\code{gs\_design.md}（设计文档）、
  \code{ideaux\_*.gp} / \code{pari\_input.gp} / \code{tlu.gp}（PARI 参考实现
  的判定实验）。
\end{itemize}
\code{supplementary/} 收录主线之外的完整求解器族（十四族相位求解器）与带日期
战役日志；\code{docs/shuttle.md} 十八次补充纪要与日志逐条对应。秘密材料与
大规模签名数据（.bin/.npy 工件，约 30\,GB）不入库；数据由 harness 对官方
提交包实现再生成（官方工件已与审计副本逐字节核对一致）。fplll 补丁版构建链与远端环境
注意事项见 \code{docs/shuttle.md} 第十四、十五次补充的工程记录。全部实验带日期戳，
日志（*.log）与纪要逐条对应。

\section{边界与不主张}

不推翻 G\"artner 2024/2052 与 Shuttle 2026/1991 的任何定理（两者的数学与伪代码经我们
独立重推与原文核对均正确）；不主张已实现等价密钥恢复（缺口已定界为 gap-1.01 的
dim-256 扭曲 uSVP 或提升方程，均超出当前算力）；GS 发散判定基于尺度递推恒等式与
真值验证的链运行（层 0/1 精确通过、$w^2$ 律与实测吻合）——是数学性质而非实现缺陷的
推诿；区分器与泄漏不直接构成 UF-CMA 伪造。对提交物的全部否定性结论（\S2、\S10 前两
行）独立成立，不依赖文献判断。设计层未破的判定（\S9.2）限定于我们已测的信道与已执行
的攻击面，不构成对设计本身的安全性证明。

\section*{参考文献}
\begin{enumerate}[label={[\arabic*]}]
\item L.~G\"artner. \emph{Compact Lattice Signatures via Iterative Rejection Sampling}.
  CRYPTO 2025（ePrint 2024/2052）.
\item Shuttle 设计论文. ePrint 2026/1991（ASIACRYPT 2026 版）.
\item L.~Ducas, A.~Durmus, T.~Lepoint, V.~Lyubashevsky. \emph{Lattice Signatures and
  Bimodal Gaussians}. CRYPTO 2013.
\item C.~Gentry, M.~Szydlo. \emph{Cryptanalysis of the Revised NTRU Signature Scheme}.
  EUROCRYPT 2002（GS 算法）.
\item A.~Lenstra, M.~Silverberg. \emph{Reconstruction of a Number from Some Values of
  its Functions}. ePrint 2014/430（L--S 张量代数路线与唯一性理论）.
\item H.~Cohn, N.~Heninger. \emph{Approximate Common Divisors via Lattices}. ANTS 2012.
\item P.-M.~Pellet-Mary, T.~Espitau 等 PARI \code{ideaux} 参考实现（GS 实现参照
  与精度判定的对照物）.
\item BLISS 泄漏家族：Flush-Gauss-Reload [GHLB+16]、Espitau 等 [EFGT17]、Pessl [Pes16]
  ——本发现的信道从微架构侧信道迁移到签名文本，为同族新成员.
\item Shuttle 提交包（NICCS 首轮）规范书与全部实现源码.
\end{enumerate}

\end{document}
